\begin{figure}[htbp]
  \centering
  \includegraphics[width=0.84\linewidth]{figures/OCRA-12x30-cut/patterns.png}
  \caption{Zbiór \texttt{OCRA-12x30-cut}: 7 wzorców o rozmiarze $12\times30$ ($n = 360$). Czarny piksel to stan $1$, biały $-1$.}
  \label{fig:OCRA-12x30-cut:patterns}
\end{figure}

\begin{figure}[htbp]
  \centering
  \includegraphics[width=\linewidth]{figures/OCRA-12x30-cut/noise_examples.png}
  \caption{Przykłady uszkodzonych wejść utworzonych z wzorca 0 zbioru \texttt{OCRA-12x30-cut}: wszystkie rodzaje i poziomy uszkodzenia.}
  \label{fig:OCRA-12x30-cut:noise}
\end{figure}

\begin{figure}[htbp]
  \centering
  \includegraphics[width=0.72\linewidth]{figures/animals-14x9/patterns.png}
  \caption{Zbiór \texttt{animals-14x9}: 6 wzorców o rozmiarze $14\times9$ ($n = 126$). Czarny piksel to stan $1$, biały $-1$.}
  \label{fig:animals-14x9:patterns}
\end{figure}

\begin{figure}[htbp]
  \centering
  \includegraphics[width=\linewidth]{figures/animals-14x9/noise_examples.png}
  \caption{Przykłady uszkodzonych wejść utworzonych z wzorca 0 zbioru \texttt{animals-14x9}: wszystkie rodzaje i poziomy uszkodzenia.}
  \label{fig:animals-14x9:noise}
\end{figure}

\begin{figure}[htbp]
  \centering
  \includegraphics[width=0.72\linewidth]{figures/large-25x25/patterns.png}
  \caption{Zbiór \texttt{large-25x25}: 6 wzorców o rozmiarze $25\times25$ ($n = 625$). Czarny piksel to stan $1$, biały $-1$.}
  \label{fig:large-25x25:patterns}
\end{figure}

\begin{figure}[htbp]
  \centering
  \includegraphics[width=\linewidth]{figures/large-25x25/noise_examples.png}
  \caption{Przykłady uszkodzonych wejść utworzonych z wzorca 0 zbioru \texttt{large-25x25}: wszystkie rodzaje i poziomy uszkodzenia.}
  \label{fig:large-25x25:noise}
\end{figure}

\begin{figure}[htbp]
  \centering
  \includegraphics[width=0.96\linewidth]{figures/large-25x25.plus/patterns.png}
  \caption{Zbiór \texttt{large-25x25.plus}: 15 wzorców o rozmiarze $25\times25$ ($n = 625$). Czarny piksel to stan $1$, biały $-1$.}
  \label{fig:large-25x25.plus:patterns}
\end{figure}

\begin{figure}[htbp]
  \centering
  \includegraphics[width=\linewidth]{figures/large-25x25.plus/noise_examples.png}
  \caption{Przykłady uszkodzonych wejść utworzonych z wzorca 0 zbioru \texttt{large-25x25.plus}: wszystkie rodzaje i poziomy uszkodzenia.}
  \label{fig:large-25x25.plus:noise}
\end{figure}

\begin{figure}[htbp]
  \centering
  \includegraphics[width=0.72\linewidth]{figures/large-25x50/patterns.png}
  \caption{Zbiór \texttt{large-25x50}: 6 wzorców o rozmiarze $25\times50$ ($n = 1250$). Czarny piksel to stan $1$, biały $-1$.}
  \label{fig:large-25x50:patterns}
\end{figure}

\begin{figure}[htbp]
  \centering
  \includegraphics[width=\linewidth]{figures/large-25x50/noise_examples.png}
  \caption{Przykłady uszkodzonych wejść utworzonych z wzorca 0 zbioru \texttt{large-25x50}: wszystkie rodzaje i poziomy uszkodzenia.}
  \label{fig:large-25x50:noise}
\end{figure}

\begin{figure}[htbp]
  \centering
  \includegraphics[width=0.96\linewidth]{figures/letters-14x20/patterns.png}
  \caption{Zbiór \texttt{letters-14x20}: 26 wzorców o rozmiarze $14\times20$ ($n = 280$). Czarny piksel to stan $1$, biały $-1$.}
  \label{fig:letters-14x20:patterns}
\end{figure}

\begin{figure}[htbp]
  \centering
  \includegraphics[width=\linewidth]{figures/letters-14x20/noise_examples.png}
  \caption{Przykłady uszkodzonych wejść utworzonych z wzorca 0 zbioru \texttt{letters-14x20}: wszystkie rodzaje i poziomy uszkodzenia.}
  \label{fig:letters-14x20:noise}
\end{figure}

\begin{figure}[htbp]
  \centering
  \includegraphics[width=0.36\linewidth]{figures/letters-abc-8x12/patterns.png}
  \caption{Zbiór \texttt{letters-abc-8x12}: 3 wzorców o rozmiarze $8\times12$ ($n = 96$). Czarny piksel to stan $1$, biały $-1$.}
  \label{fig:letters-abc-8x12:patterns}
\end{figure}

\begin{figure}[htbp]
  \centering
  \includegraphics[width=\linewidth]{figures/letters-abc-8x12/noise_examples.png}
  \caption{Przykłady uszkodzonych wejść utworzonych z wzorca 0 zbioru \texttt{letters-abc-8x12}: wszystkie rodzaje i poziomy uszkodzenia.}
  \label{fig:letters-abc-8x12:noise}
\end{figure}

\begin{figure}[htbp]
  \centering
  \includegraphics[width=0.60\linewidth]{figures/small-7x7/patterns.png}
  \caption{Zbiór \texttt{small-7x7}: 5 wzorców o rozmiarze $7\times7$ ($n = 49$). Czarny piksel to stan $1$, biały $-1$.}
  \label{fig:small-7x7:patterns}
\end{figure}

\begin{figure}[htbp]
  \centering
  \includegraphics[width=\linewidth]{figures/small-7x7/noise_examples.png}
  \caption{Przykłady uszkodzonych wejść utworzonych z wzorca 0 zbioru \texttt{small-7x7}: wszystkie rodzaje i poziomy uszkodzenia.}
  \label{fig:small-7x7:noise}
\end{figure}
